\documentclass[11pt]{article}

% Shared notes settings for Elements of AIML lectures (UPES).
% Keep this file free of \documentclass to allow reuse via \input.

\usepackage[utf8]{inputenc}
\usepackage[T1]{fontenc}
\usepackage[a4paper,margin=1in]{geometry}
\usepackage{hyperref}
\usepackage{xurl}
\usepackage{booktabs}
\usepackage{amsmath}
\usepackage{amssymb}
\usepackage{listings}
\usepackage{xcolor}
\usepackage{enumitem}
\usepackage{parskip}

% Code listings (general-purpose).
\lstset{
  basicstyle=\ttfamily\small,
  keywordstyle=\color{blue},
  commentstyle=\color{gray},
  stringstyle=\color{teal},
  showstringspaces=false,
  columns=fullflexible,
  frame=single,
  framerule=0.2pt,
  breaklines=true
}

\setlist[itemize]{topsep=0.3em,itemsep=0.2em}
\setlist[enumerate]{topsep=0.3em,itemsep=0.2em}

% Simple per-section source line.
\newcommand{\notesources}[1]{\par\noindent{\small\color{gray}\textit{Sources:} \sloppy #1\par}}

% Unit-wise source strings for this subject (used by slides/notes).
% Path is relative to the lecture `latex/` folder (the compilation CWD).
\IfFileExists{../../../_shared/aiml-sources.tex}{%
  % Source strings for Elements of AIML (used by slides + notes).

\newcommand{\AIMLRepoURL}{https://github.com/tali7c/Elements-of-AIML}

\newcommand{\AIMLLabSources}{%
Provost and Fawcett, \textit{Data Science for Business}.;%
McKinney, \textit{Python for Data Analysis}.;%
WEKA Documentation (Waikato Environment for Knowledge Analysis), accessed 2026-02-28.;%
Matplotlib and Seaborn documentation, accessed 2026-02-28.%
}

\newcommand{\AIMLUnitISources}{%
Rich and Knight, \textit{Artificial Intelligence}, McGraw Hill, 2017.;%
Russell and Norvig, \textit{Artificial Intelligence: A Modern Approach} (4e), Ch. 1--2 (Introduction, Intelligent Agents).;%
Poole and Mackworth, \textit{Artificial Intelligence: Foundations of Computational Agents} (2e), selected sections.%
}

\newcommand{\AIMLUnitIISources}{%
Russell and Norvig, \textit{Artificial Intelligence: A Modern Approach} (4e), Ch. 7--9 (Logical Agents, First-Order Logic, Inference).;%
Huth and Ryan, \textit{Logic in Computer Science} (2e), selected sections on propositional and first-order logic.;%
Genesereth and Nilsson, \textit{Logical Foundations of Artificial Intelligence}, selected chapters.%
}

\newcommand{\AIMLUnitIIISources}{%
Mueller and Massaron, \textit{Machine Learning for Dummies}, For Dummies, 2016.;%
Mitchell, \textit{Machine Learning}, selected chapters on learning basics and evaluation.;%
Scikit-learn documentation, accessed 2026-02-28.%
}

\newcommand{\AIMLUnitIVSources}{%
Mueller and Massaron, \textit{Machine Learning for Dummies}, For Dummies, 2016.;%
James, Witten, Hastie, and Tibshirani, \textit{An Introduction to Statistical Learning} (2e), selected chapters.;%
Scikit-learn documentation, accessed 2026-02-28.%
}

\newcommand{\AIMLUnitVSources}{%
NIST, \textit{AI Risk Management Framework (AI RMF 1.0)}, 2023.;%
Barocas, Hardt, and Narayanan, \textit{Fairness and Machine Learning} (online book), selected sections.;%
Knaflic, \textit{Storytelling with Data}, selected chapters on communicating results.%
}
%
}{%
  % Faculty copies live under `private/` so their compile CWD is different.
  \IfFileExists{../../../../public/_shared/aiml-sources.tex}{%
    % Source strings for Elements of AIML (used by slides + notes).

\newcommand{\AIMLRepoURL}{https://github.com/tali7c/Elements-of-AIML}

\newcommand{\AIMLLabSources}{%
Provost and Fawcett, \textit{Data Science for Business}.;%
McKinney, \textit{Python for Data Analysis}.;%
WEKA Documentation (Waikato Environment for Knowledge Analysis), accessed 2026-02-28.;%
Matplotlib and Seaborn documentation, accessed 2026-02-28.%
}

\newcommand{\AIMLUnitISources}{%
Rich and Knight, \textit{Artificial Intelligence}, McGraw Hill, 2017.;%
Russell and Norvig, \textit{Artificial Intelligence: A Modern Approach} (4e), Ch. 1--2 (Introduction, Intelligent Agents).;%
Poole and Mackworth, \textit{Artificial Intelligence: Foundations of Computational Agents} (2e), selected sections.%
}

\newcommand{\AIMLUnitIISources}{%
Russell and Norvig, \textit{Artificial Intelligence: A Modern Approach} (4e), Ch. 7--9 (Logical Agents, First-Order Logic, Inference).;%
Huth and Ryan, \textit{Logic in Computer Science} (2e), selected sections on propositional and first-order logic.;%
Genesereth and Nilsson, \textit{Logical Foundations of Artificial Intelligence}, selected chapters.%
}

\newcommand{\AIMLUnitIIISources}{%
Mueller and Massaron, \textit{Machine Learning for Dummies}, For Dummies, 2016.;%
Mitchell, \textit{Machine Learning}, selected chapters on learning basics and evaluation.;%
Scikit-learn documentation, accessed 2026-02-28.%
}

\newcommand{\AIMLUnitIVSources}{%
Mueller and Massaron, \textit{Machine Learning for Dummies}, For Dummies, 2016.;%
James, Witten, Hastie, and Tibshirani, \textit{An Introduction to Statistical Learning} (2e), selected chapters.;%
Scikit-learn documentation, accessed 2026-02-28.%
}

\newcommand{\AIMLUnitVSources}{%
NIST, \textit{AI Risk Management Framework (AI RMF 1.0)}, 2023.;%
Barocas, Hardt, and Narayanan, \textit{Fairness and Machine Learning} (online book), selected sections.;%
Knaflic, \textit{Storytelling with Data}, selected chapters on communicating results.%
}
%
  }{%
    % Source strings for Elements of AIML (used by slides + notes).

\newcommand{\AIMLRepoURL}{https://github.com/tali7c/Elements-of-AIML}

\newcommand{\AIMLLabSources}{%
Provost and Fawcett, \textit{Data Science for Business}.;%
McKinney, \textit{Python for Data Analysis}.;%
WEKA Documentation (Waikato Environment for Knowledge Analysis), accessed 2026-02-28.;%
Matplotlib and Seaborn documentation, accessed 2026-02-28.%
}

\newcommand{\AIMLUnitISources}{%
Rich and Knight, \textit{Artificial Intelligence}, McGraw Hill, 2017.;%
Russell and Norvig, \textit{Artificial Intelligence: A Modern Approach} (4e), Ch. 1--2 (Introduction, Intelligent Agents).;%
Poole and Mackworth, \textit{Artificial Intelligence: Foundations of Computational Agents} (2e), selected sections.%
}

\newcommand{\AIMLUnitIISources}{%
Russell and Norvig, \textit{Artificial Intelligence: A Modern Approach} (4e), Ch. 7--9 (Logical Agents, First-Order Logic, Inference).;%
Huth and Ryan, \textit{Logic in Computer Science} (2e), selected sections on propositional and first-order logic.;%
Genesereth and Nilsson, \textit{Logical Foundations of Artificial Intelligence}, selected chapters.%
}

\newcommand{\AIMLUnitIIISources}{%
Mueller and Massaron, \textit{Machine Learning for Dummies}, For Dummies, 2016.;%
Mitchell, \textit{Machine Learning}, selected chapters on learning basics and evaluation.;%
Scikit-learn documentation, accessed 2026-02-28.%
}

\newcommand{\AIMLUnitIVSources}{%
Mueller and Massaron, \textit{Machine Learning for Dummies}, For Dummies, 2016.;%
James, Witten, Hastie, and Tibshirani, \textit{An Introduction to Statistical Learning} (2e), selected chapters.;%
Scikit-learn documentation, accessed 2026-02-28.%
}

\newcommand{\AIMLUnitVSources}{%
NIST, \textit{AI Risk Management Framework (AI RMF 1.0)}, 2023.;%
Barocas, Hardt, and Narayanan, \textit{Fairness and Machine Learning} (online book), selected sections.;%
Knaflic, \textit{Storytelling with Data}, selected chapters on communicating results.%
}
%
  }%
}


\title{Elements of AIML\\Unit 05 -- Lecture 01 Notes\\AI for Society, Women \& Environment}
\author{Tofik Ali}
\date{\today}

\begin{document}
\maketitle
\tableofcontents

\section{Lecture Overview}
This lecture introduces responsible AI considerations:
\begin{itemize}[leftmargin=*]
  \item fairness and bias,
  \item privacy and security,
  \item safety and accountability,
  \item environmental sustainability,
  \item inclusive design considerations (including gender impacts).
\end{itemize}
\notesources{\AIMLUnitVSources}

\section{How to use these notes (self-study)}
Responsible AI is best studied through questions:
Who is affected, what harms are possible, and what evidence do we need before deployment?
When you read each section (fairness, privacy, sustainability), try to connect it to a concrete scenario such as loan approval, hiring, or healthcare triage.

\section{Before you start}
These notes assume you already understand basic model evaluation from Unit 03, especially precision, recall, false positives, false negatives, and why metrics depend on the application.
If those ideas feel weak, briefly revisit the evaluation material in Unit 03 before reading this lecture in depth.

\section{Case-study lens used in this lecture}
Use the same six questions throughout the lecture:
\begin{itemize}[leftmargin=*]
  \item \textbf{Task}: what decision or prediction is the model making?
  \item \textbf{Data}: what inputs and labels are used?
  \item \textbf{Metric}: what performance measure matters most?
  \item \textbf{Main risk}: what harm can happen if the system fails?
  \item \textbf{Monitoring signal}: what should we track after deployment?
  \item \textbf{Human role}: where should review, appeal, or override exist?
\end{itemize}

\paragraph{Case-study snapshot: loan screening}
\textbf{Task:} predict whether an applicant is likely to repay a loan.\\
\textbf{Data:} income history, debt ratio, repayment history, employment pattern.\\
\textbf{Metric:} recall for risky applicants and false positive rate for rejected safe applicants.\\
\textbf{Main risk:} unfair rejection of capable applicants from under-represented groups.\\
\textbf{Monitoring signal:} approval rate, default rate, and false positive rate by subgroup over time.\\
\textbf{Human role:} borderline cases should go to manual review with documented justification.

\paragraph{Case-study snapshot: hospital triage support}
\textbf{Task:} flag patients who may need urgent specialist attention.\\
\textbf{Data:} symptoms, age, test results, medical history.\\
\textbf{Metric:} recall/sensitivity, because missed urgent cases are costly.\\
\textbf{Main risk:} false negatives can delay treatment and cause harm.\\
\textbf{Monitoring signal:} missed-critical-case rate, subgroup recall, and drift in incoming patient profiles.\\
\textbf{Human role:} the model should support, not replace, clinician judgment.

\section{Learning Outcomes}
\begin{itemize}[leftmargin=*]
  \item Explain why responsible AI is required beyond accuracy metrics.
  \item Identify sources of bias and describe fairness at a high level.
  \item Describe privacy risks in data-driven systems.
  \item Explain environmental trade-offs of training and deploying models.
\end{itemize}

\section{Keywords}
\begin{itemize}[leftmargin=*]
  \item \textbf{Bias:} systematic error or unfairness introduced by data/processes.
  \item \textbf{Fairness:} notion of equitable treatment across individuals or groups.
  \item \textbf{Privacy:} protection of sensitive information about individuals.
  \item \textbf{Accountability:} ability to assign responsibility, explain and audit decisions.
  \item \textbf{Sustainability:} considering energy/compute and environmental costs.
\end{itemize}

\section{Abbreviations and key terms}
\begin{itemize}[leftmargin=*]
  \item \textbf{PII}: Personally Identifiable Information (e.g., phone number, address, ID number).
  \item \textbf{Drift}: data distribution changes over time (inputs or labels).
  \item \textbf{Audit}: systematic check of datasets/models for issues (bias, leakage, failures).
  \item \textbf{Stakeholders}: people affected by a system (users, applicants, patients, society).
  \item \textbf{High-stakes}: settings where errors can cause significant harm (health, finance, justice).
\end{itemize}

\section{Why responsible AI?}
AI systems influence human lives. In many domains, errors have different costs for different people.
Therefore we must evaluate:
\begin{itemize}[leftmargin=*]
  \item performance (accuracy/error),
  \item fairness across relevant groups,
  \item privacy and data governance,
  \item safety and security,
  \item monitoring and accountability.
\end{itemize}

\subsection{A practical mindset}
Responsible AI is not a single technique; it is a \textbf{process}:
\begin{itemize}[leftmargin=*]
  \item define harms and success criteria with stakeholders,
  \item measure model performance \emph{overall and by subgroup},
  \item implement mitigations (data, model, thresholds, policy),
  \item monitor and update after deployment.
\end{itemize}

\section{Bias and fairness}
\subsection{Sources of bias}
Common bias sources include:
\begin{itemize}[leftmargin=*]
  \item \textbf{Sampling bias}: training data under-represents some groups.
  \item \textbf{Measurement/label bias}: labels reflect human errors or biased judgments.
  \item \textbf{Historical bias}: past patterns embed unfairness into data.
  \item \textbf{Proxy bias}: non-sensitive features correlate with sensitive attributes.
\end{itemize}

\subsection{Fairness (practical view)}
Fairness is context-dependent. Common questions:
\begin{itemize}[leftmargin=*]
  \item Are error rates similar across groups?
  \item Do different groups receive systematically different outcomes?
  \item Is there a strong justification for differences?
\end{itemize}
There is no single metric that solves every fairness situation; we choose based on domain and risk tolerance.
In the loan-screening example above, equal overall accuracy is not enough if one subgroup experiences a much higher false rejection rate.

\subsection{Fairness metrics (high-level examples)}
Some common (but not universal) fairness notions:
\begin{itemize}[leftmargin=*]
  \item \textbf{Demographic parity:} positive prediction rates are similar across groups.
  \item \textbf{Equal opportunity:} true positive rates (recall) are similar across groups.
  \item \textbf{Error rate balance:} false positive/false negative rates are similar across groups.
\end{itemize}
Important: different metrics can conflict; you must choose based on domain harms, legal constraints, and feasibility.

\subsection{Mitigation options (practical)}
Examples of fairness-related mitigations:
\begin{itemize}[leftmargin=*]
  \item \textbf{Data fixes}: collect more representative data, reweight samples, improve labels.
  \item \textbf{Modeling fixes}: add constraints/regularization, use interpretable models where needed.
  \item \textbf{Decision fixes}: different thresholds by group (requires strong justification), human review for borderline cases.
  \item \textbf{Process fixes}: documentation, audits, and continuous monitoring.
\end{itemize}

\section{Privacy and security}
\subsection{Dataset privacy}
Sensitive data (medical, financial) must be handled with governance:
\begin{itemize}[leftmargin=*]
  \item access controls, consent, retention policies,
  \item encryption and secure storage,
  \item careful sharing with third parties.
\end{itemize}

\subsection{Privacy vs security (quick distinction)}
\begin{itemize}[leftmargin=*]
  \item \textbf{Privacy}: who is allowed to know what about an individual.
  \item \textbf{Security}: preventing unauthorized access, tampering, and attacks.
\end{itemize}
Both are needed for real-world AI systems.

\subsection{Model privacy (important)}
Even if the dataset is protected, models can leak information:
\begin{itemize}[leftmargin=*]
  \item membership inference (whether a person was in training data),
  \item memorization of rare records,
  \item extraction attacks when models are exposed publicly.
\end{itemize}
Mitigations include limiting access, careful logging, and privacy-preserving techniques (advanced topic).
In the hospital triage example, privacy is not optional because even a small leak can expose highly sensitive medical details.

\section{Safety and accountability}
In high-stakes settings:
\begin{itemize}[leftmargin=*]
  \item humans may need to remain in the loop,
  \item decisions should be explainable and contestable,
  \item monitoring should detect drift and harmful outcomes.
\end{itemize}

\subsection{Accountability artifacts (what to keep)}
In practice, organizations maintain:
\begin{itemize}[leftmargin=*]
  \item dataset documentation (source, time range, known limitations),
  \item model documentation (features used, training procedure, evaluation results),
  \item decision logs (for audit/appeals), where appropriate and legal.
\end{itemize}

\section{Environmental impact and sustainability}
Compute has real energy costs. Practical ways to reduce impact:
\begin{itemize}[leftmargin=*]
  \item prefer smaller models when sufficient,
  \item reduce unnecessary features and retraining,
  \item use efficient evaluation and early stopping where appropriate,
  \item measure cost vs benefit (marginal accuracy gains vs compute).
\end{itemize}

\subsection{Sustainability is also a design choice}
Sometimes the most responsible choice is not the most complex model:
\begin{itemize}[leftmargin=*]
  \item simpler models can be easier to audit and explain,
  \item smaller compute footprints reduce cost and environmental impact,
  \item marginal accuracy gains may not justify the added cost/risk.
\end{itemize}

\section{Inclusive design (women \& society)}
Under-representation can cause unequal performance. Practical actions:
\begin{itemize}[leftmargin=*]
  \item check representation in datasets,
  \item evaluate metrics by subgroup,
  \item collect feedback and involve stakeholders in defining success.
\end{itemize}
This is why the monitoring signals listed in the case-study lens should always be checked by subgroup, not only on the full population average.

\section{Practice (with answers)}
\subsection*{P1. Give one example of a proxy feature that could introduce unfairness.}
\textbf{Answer:} ZIP code can act as a proxy for socioeconomic status or race/ethnicity in some regions; using it for credit decisions can introduce unfair outcomes.

\subsection*{P2. What is drift and why does it matter?}
\textbf{Answer:} drift is change in input/label distribution over time; it can reduce model accuracy and increase unfair outcomes if not monitored.

\section{Common pitfalls (and how to avoid them)}
\begin{itemize}[leftmargin=*]
  \item \textbf{Treating fairness as one metric:} fairness is context-dependent; different metrics can conflict.
  \item \textbf{Assuming anonymization is always safe:} re-identification is possible; privacy needs governance and access controls.
  \item \textbf{Ignoring subgroup performance:} a model can be ``good on average'' but harmful for some groups.
  \item \textbf{No monitoring plan:} drift can turn a safe model into an unsafe one after deployment.
  \item \textbf{Over-optimizing accuracy:} marginal accuracy gains may increase cost, risk, and environmental impact.
\end{itemize}

\section{Exit questions (with answers)}

\subsection*{Q1. List three sources of bias in ML datasets.}
\textbf{Answer:} sampling bias, measurement/label bias, historical bias, proxy bias (any three).

\subsection*{Q2. Give one example where fairness is critical.}
\textbf{Answer:} loan approval, hiring screening, medical triage, criminal justice risk assessment.

\subsection*{Q3. One privacy risk of models (not just datasets).}
\textbf{Answer:} membership inference or memorization: the model can leak information about specific training examples.

\subsection*{Q4. Two ways to reduce environmental impact.}
\textbf{Answer:} use smaller models, reduce retraining frequency, optimize pipelines, prefer efficient algorithms/hardware, avoid unnecessary experiments.

\section{Summary}
Responsible AI means evaluating and managing fairness, privacy, safety, and sustainability along with accuracy.
\notesources{\AIMLUnitVSources}

\end{document}
